\chapter{Einleitung}

Wir leben in einer Zeit der Digitalisierung. Immer mehr Bereiche des Lebens werden durch digitale Lösungen vorangetrieben.
Dazu gehören auch Lernveranstaltungen. Mithilfe von Webseiten wie Moodle können Vorlesungsinhalte sowie Testate digital angeboten werden.
Besonders in Fachbereichen wie der Informatik bietet es sich an, Vorlesungen und Klausuren zu digitalisieren, da der Code schließlich digital und nicht auf Papier ausgeführt werden soll.
Daher soll für den Kurs „Programmiertechnik 1” zukünftig die Lernplattform Codetask verwendet werden. 
Mit Codetask sollen Vorlesungen und Klausuren durchgeführt werden. Der Schwerpunkt von Codetask liegt auf der Programmiersprache Scala.
In den Vorlesungen werden Inhalte vorgetragen und passende Aufgaben zum Thema gestellt, die von den Studierenden bearbeitet werden sollen.
Zu diesen Aufgaben gehören Single-Choice-Fragen, Multi-Choice-Fragen, Textaufgaben und Lückentextaufgaben, die als „Koans” bezeichnet werden.
Eine weitere Aufgabenart, die bei Codetask eine Besonderheit ist, sind Codeaufgaben, bei denen die Studierenden einen funktionierenden Code schreiben müssen.
Dieser Code wird dann durch einen internen Compiler ausgewertet, um zu prüfen, ob er funktioniert. Diese Art der Aufgabe wird „Codetask” genannt.
Die gleichen Aufgabenarten sollen auch in den Klausuren gestellt werden.

Da Vorlesungen und insbesondere Klausuren mit Codetask durchgeführt werden sollen, muss Codetask auf einem zuverlässigen System betrieben werden.
Eine Klausur ist ein kritischer Bereich, in dem die Anwendung verfügbar bleiben muss.
Die Anwendung muss dabei stabil gehalten werden, damit es bei einer Klausur zu keinen Ausfällen kommt. Falls doch welche auftreten, muss die Anwendung schnell wiederhergestellt werden.
Fehler im System dürfen die Klausuren nicht beeinflussen.

Aktuell wird Codetask auf einem Docker Compose betrieben.
Ein Docker Compose bietet jedoch keine optimale Ausfallsicherheit, da jeder Service nur einmal ausgeführt wird und die gesamte Anwendung nur auf einem Server läuft.
Dadurch ist die Last der Anwendung auf einen Server ausgelegt und an dessen Grenzen limitiert.
Dabei ist inzwischen schon normalität, Anwendungen in verteilten Systemen zu betreiben. 
Der benötigte Ressourcenverbrauch erhöht sich schneller in der Gesellschaft, als die Hardwareentwicklung sich an diesen Ressourcenverbrauch anpassen kann.
Es ist aber oft auch günstiger, mehrere Hardware-Systeme mit normalem Ressourcenverbrauch zu beschaffen, als ein System, das einen größeren Ressourcenverbrauch verwalten kann.
Seine Anwendung auf mehrere Systeme zu verteilen, bringt aber nicht nur die Lastverteilung, sondern macht auch die Anwendung ausfallsicherer.
Fällt eine Hardware eines Systems aus, kann die Anwendung auf die andere Hardware weiterbetrieben werden.
Dadurch stellt sich die Frage, ob Codetask für die Produktionsumgebung in einem verteilten System besser aufgehoben ist.
Um Codetask in einem verteilten System zu betreiben, soll Kubernetes verwendet werden. Kubernetes selbst bietet dabei auch viele Features, mit denen die Zuverlässigkeit von Codetask zunehmen kann.

Daher beschäftigt sich diese Arbeit mit der Umstellung von Codetask von einer Docker-Compose-Umgebung auf eine funktionierende Kubernetes-Umgebung.
Zu diesem Zweck muss zunächst ein Kubernetes-Cluster im Hochschulnetzwerk ohne Cloud-Einbindung erstellt werden. Dieser soll so eingerichtet werden, dass eine höhere Zuverlässigkeit erzielt wird.
Abschließend sollen die Docker-Compose-Umgebung und die Kubernetes-Umgebung durch Tests verglichen werden. Dabei wird evaluiert, wie sich die Zuverlässigkeit in den Kategorien Verfügbarkeit, Fehlertoleranz, Stabilität und Wiederherstellbarkeit unterscheidet.

Zunächst werden die wichtigsten Grundlagen von Docker Compose und Kubernetes sowie die Faktoren, die die Zuverlässigkeit beeinflussen, vorgestellt. Anschließend wird erläutert, wie Kubernetes die Zuverlässigkeit verbessern kann.
Anschließend wird das aktuelle System, das mit Docker Compose betrieben wird, analysiert und vorgestellt.
Daraufhin wird die Implementierung der Kubernetes-Umgebung vorgestellt und dabei auch erörtert, wie sich dadurch die Zuverlässigkeit erhöhen kann.
Nach der Implementierung der Kubernetes-Umgebung werden beide Umgebungen mit Load- und Chaostests getestet und die Ergebnisse verglichen.
Zum Schluss folgt das Fazit, in dem die Einschätzung präsentiert wird, ob sich die Umstellung rentiert, und aufgezeigt wird, wie darauf aufgebaut werden kann.
